\section{The Architect's Question}\label{the-architects-question}

Before we can place a model across multiple GPUs or across a cluster, we
need to know what the interconnect can --- and cannot --- move. This
chapter asks: given a model size, a parallelism strategy, and a traffic
pattern (all-reduce, all-gather, reduce-scatter), what is the
communication time, and does it dominate the compute budget? We will
size the hidden bottleneck that often goes unnoticed until the job
stalls.

\section{Concept}\label{concept}

Communication in a distributed LLM system is not a single pipe. It is a
hierarchy: each GPU talks to its neighbors over PCIe, nodes within a
server meet over NVLink/NVSwitch, racks talk over 200/400 GbE RoCE, and
cross-rack or cross-data-center traffic rides InfiniBand. Each layer has
a distinct bandwidth ceiling, a distinct latency profile, and a distinct
cost per gigabyte. The architect must match the collective operation to
the layer where the bandwidth is sufficient, because moving the same
data over a slower layer adds minutes to training time or inference
latency for no computational benefit.

The collective operations that dominate multi-node training and
inference are:

\begin{itemize}
\tightlist
\item
  \textbf{all-reduce}: every GPU contributes a gradient fragment; the
  result (the sum, scaled by 1/N) is delivered to every GPU. This is the
  workhorse of data-parallel and model-parallel training.
\item
  \textbf{all-gather}: every GPU contributes a tensor; every GPU
  receives the concatenation of all fragments. Used when a model split
  across nodes needs the full weight or KV cache locally.
\item
  \textbf{reduce-scatter}: the inverse of all-gather. Every GPU
  contributes a fragment; each GPU receives a portion of the final
  result, sized 1/N of the input. Used when the full tensor does not
  need to reside on any single GPU.
\end{itemize}

The bandwidth each collective consumes depends on three factors: the
data volume (bytes), the number of participating nodes (N), and the
interconnect's effective bandwidth in the presence of contention. A 70B
FP16 model is 140 GB. An all-reduce of that model across 8 GPUs on a
single node moves 140 GB over NVLink; across 8 nodes connected by
InfiniBand moves 140 GB (or more, depending on the algorithm) over the
fabric. The same 140 GB over a 25 GB/s Ethernet link takes more than 5×
the time of the same operation over a 900 GB/s NVLink.

\textbf{\hyperref[tab:9.1]{Table 9-1}} --- Interconnect hierarchy with bandwidth
(bidirectional where applicable), latency, and approximate cost per GB.
PCIe Gen5 x16 \textasciitilde64 GB/s, \textasciitilde1 µs latency, low
cost (integrated); NVLink H100 900 GB/s bisection, \textasciitilde0.5--1
µs latency, high cost (GPU-attached); Ethernet RoCE2 200/400 Gb/s ≈
25/50 GB/s, \textasciitilde10--25 µs latency, moderate cost (standard
NICs); InfiniBand HDR 400 Gb/s / NDR 800 Gb/s ≈ 50/100 GB/s,
\textasciitilde2--5 µs latency, high cost (switches, optics). Cost per
GB is inversely proportional to bandwidth for fixed-form optics, but
system cost includes cables, switches, and rack density trade-offs.

\subsubsection{\hyperref[tab:9.1]{Table 9-1} --- The interconnect
hierarchy}\label{table-9-1-the-interconnect-hierarchy}

\begin{longtable}[]{@{}
  >{\raggedright\arraybackslash}p{0.1667\linewidth}
  >{\raggedright\arraybackslash}p{0.1667\linewidth}
  >{\raggedright\arraybackslash}p{0.1667\linewidth}
  >{\raggedright\arraybackslash}p{0.1667\linewidth}
  >{\raggedright\arraybackslash}p{0.1667\linewidth}
  >{\raggedright\arraybackslash}p{0.1667\linewidth}@{}}
\toprule\noalign{}
\begin{minipage}[b]{\linewidth}\raggedright
Layer
\end{minipage} & \begin{minipage}[b]{\linewidth}\raggedright
Example
\end{minipage} & \begin{minipage}[b]{\linewidth}\raggedright
Bandwidth (one direction / aggregate)
\end{minipage} & \begin{minipage}[b]{\linewidth}\raggedright
Latency
\end{minipage} & \begin{minipage}[b]{\linewidth}\raggedright
Cost
\end{minipage} & \begin{minipage}[b]{\linewidth}\raggedright
{[}source{]}
\end{minipage} \\
\midrule\noalign{}
\endhead
\bottomrule\noalign{}
\endlastfoot
Chip-to-memory & HBM3 (H100) & 3.35 TB/s & \textasciitilde ns & --- &
{[}2° FACT{]} \\
GPU-to-GPU (in-node) & NVLink (H100) & 900 GB/s per GPU, bidirectional
aggregate & \textasciitilde us & high capex & {[}2° FACT{]} \\
GPU-to-CPU / NIC & PCIe Gen5 & \textasciitilde64 GB/s per x16 direction
& \textasciitilde us & low & {[}2° FACT{]} \\
Node-to-node (rack) & RoCE / Ethernet 400Gb/s & \textasciitilde50 GB/s
per port & \textasciitilde us & medium & {[}2° FACT{]} \\
Rack/cluster & InfiniBand HDR/NDR & 400/800 Gb/s ≈ 50/100 GB/s per port
& \textasciitilde us & high capex & {[}2° FACT{]} \\

\label{tab:9.1}\end{longtable}

\emph{(Hierarchy: bandwidth drops \textasciitilde2--3 orders of
magnitude from HBM to the cluster fabric --- the reason the
interconnect, not the GPU, often binds at scale. Figures are established
hardware specs {[}2° FACT{]}.)}

\section{Mental Model}\label{mental-model}

Think of the interconnect as a series of concentric rings. The GPU sits
at the center, reachable fastest via PCIe. One ring out is the node:
GPUs on the same server are joined by NVLink and NVSwitch, forming a
high-radix, low-latency mesh. The next ring out is the rack: 10--20 GbE
or 100/200/400 GbE RoCE connects nodes. The outermost ring is the
cluster or data centre: InfiniBand HDR (400 Gb/s) or NDR (800 Gb/s)
stitches racks into a single logical network. A well-designed deployment
keeps the heavy collectives (all-reduce, reduce-scatter) on the inner
rings and reserves the outer rings for less volume or latency-tolerant
traffic.

The arithmetic is relentless. If all-reduce moves 140 GB and the NVLink
bisection provides 900 GB/s per direction, \textbf{a lower bound} on
time is 140 GB ÷ 900 GB/s ≈ \textbf{0.16 s {[}2° DERIVED{]}} --- note
the word \emph{lower bound}. The actual completion time of a real
collective is higher because a collective is not one flat transfer but a
sequence of message-passing and reduction steps across the topology:

\[
T_{\text{collective}} \approx \alpha \times n_\text{steps} + \frac{\text{bytes}}{\beta}
\]

where \(\alpha\) is the per-step latency (message setup,
synchronization), \(\beta\) is the achieved (not peak) bandwidth, and
\(n_\text{steps}\) is the number of message-passing/reduction hops. (For
a flat transfer \(n_\text{steps}=1\).) The 0.16 s figure uses only the
\(\text{bytes}/\beta\) term at peak bandwidth; real NCCL all-reduce on
NVLink lands tens of \% above it once per-step overhead, ring stages,
and topology routing are included. So treat the
\texttt{bytes\ ÷\ bandwidth} form as a \emph{sizing lower bound}, and
reserve measured \texttt{nccl-tests} numbers (Section 4) for capacity
decisions. The relative ranking between interconnects is what matters
most here: if those same 140 GB travel over a 400 Gb/s InfiniBand link
(≈ 50 GB/s effective unidirectional) it is ≈ 2.8 s --- 18× longer; over
25 Gb/s Ethernet (≈ 3 GB/s unidirectional) ≈ 47 s. The interconnect
choice is a first-order determinant of wall-clock time in all cases.

\begin{figure}
\centering
\pandocbounded{\includegraphics[keepaspectratio,alt={\hyperref[fig:9.1]{Fig 9.1} --- All-reduce time vs data volume, by interconnect tier {[}2° DERIVED{]}},width=\maxwidth{\textwidth},keepaspectratio]{figures/fig-09-0901.pdf}}
\caption{\hyperref[fig:9.1]{Fig 9.1} --- All-reduce time vs data volume, by interconnect
tier {[}2° DERIVED{]}}

\label{fig:9.1}\end{figure}

\emph{\hyperref[fig:9.1]{Fig 9.1} --- All-reduce completion time as a function of data
volume (x-axis, log GB) and interconnect effective bandwidth. The four
lines --- NVSwitch 1.8 TB/s, NVLink 0.9 TB/s, InfiniBand 0.4 TB/s,
Ethernet 0.1 TB/s --- are peak-bandwidth lower bounds (α/β caveat in
§9.3); the dashed marker sits at the 140 GB weight footprint of a 70B
model (≈ 70 GB of weights in BF16 across 8 participants). {[}2°
DERIVED{]}}

\section{Worked Example}\label{worked-example}

Consider the canonical scenario: one host eight × H100, each GPU holding
140 GB of model weights in FP16. We will run data-parallel training with
full model replicas, so every all-reduce moves the full gradient tensor
of 140 GB. Let us compute the communication time under three
interconnect options.

\textbf{Option A --- NVLink on a single host.} H100 GPUs provide 900
GB/s bidirectional NVLink bandwidth per GPU, and NVSwitch aggregates
bisection bandwidth. For an all-reduce of 140 GB across 8 GPUs on one
host, the effective bandwidth is close to the per-link peak because
NVSwitch can forward multiple concurrent channels without severe
contention. The time is approximately:

\[
T = \frac{\text{bytes}}{B_\text{eff}} = \frac{140 \text{ GB}}{900 \text{ GB/s}} \approx 0.16 \text{ s}
\]

\textbf{Option B --- InfiniBand HDR across 8 nodes, one GPU per node.}
HDR InfiniBand delivers 400 Gb/s per link ≈ 50 GB/s unidirectional after
8b/10b encoding and protocol headers. A ring-all-reduce or
tree-all-reduce across 8 nodes will see lower effective bandwidth due to
multi-hop forwarding and congestion. A realistic effective bandwidth is
roughly 40 GB/s. The time is:

\[
T = \frac{140 \text{ GB}}{40 \text{ GB/s}} \approx 3.5 \text{ s}
\]

\textbf{Option C --- 25 Gb/s Ethernet (RoCE2) across 8 nodes.} 25 Gb/s ≈
3.125 GB/s unidirectional. With protocol overhead and multi-node tree
contention, a practical effective bandwidth might be 2.5 GB/s. The time
is:

\[
T = \frac{140 \text{ GB}}{2.5 \text{ GB/s}} \approx 56 \text{ s}
\]

These numbers illustrate the scale: moving from NVLink to InfiniBand
adds \textasciitilde22× latency, and forcing the same all-reduce over
Ethernet adds \textasciitilde350× latency. In a training job where 100
all-reduce steps run per iteration, the NVLink path adds
\textasciitilde16 s of communication per iteration, the InfiniBand path
adds \textasciitilde350 s, and the Ethernet path adds
\textasciitilde5,600 s. The latter two would effectively serialize the
training, turning a minutes-long job into an hours- or days-long one.

The worked example also highlights why model parallelism and pipeline
parallelism exist: when a single GPU cannot hold the model, splitting it
across nodes moves the communication from the inner rings to the outer
rings, and the architect must budget the extra seconds per collective
against the savings from fitting a larger model.

\section{Measurement}\label{measurement}

Measuring communication in situ is harder than counting tokens, because
the bandwidth numbers above are peaks; real systems see lower effective
bandwidth due to congestion, operating-system interference, and protocol
overhead. Three practical measurements an architect can make:

\begin{enumerate}
\def\labelenumi{\arabic{enumi}.}
\item
  \textbf{Bandwidth benchmark per link.} Use \texttt{nccl-tests} or
  \texttt{ib\_write\_bw} / \texttt{bwperf} to measure achieved bandwidth
  on PCIe, NVLink (via NVIDIA's Nsight), RoCE, and InfiniBand. Record
  the unidirectional and bidirectional throughput for payload sizes
  matching the collective (typically 1 MB--1 GB). These numbers anchor
  the arithmetic to the real hardware.
\item
  \textbf{Collective latency and throughput.} Run \texttt{nccl-tests}
  all-reduce micro‑benchmarks varying the number of GPUs (1, 2, 4, 8)
  and the tensor size (from 100 MB up to 140 GB). Plot time vs.~size and
  extract the slope --- that slope is the effective bandwidth the
  collective will see in production. Compare the slope across 1 node, 2
  nodes, 4 nodes, to see where the interconnect ceiling is hit.
\item
  \textbf{End-to-end training iteration time.} Instrument the training
  loop to separate compute time (FLOPs) from communication time
  (all-reduce collectives). Many frameworks (PyTorch Distributed, vLLM,
  Deepspeed) already expose a per-iteration breakdown. The communication
  time per iteration, divided by the number of all-reduce steps, gives
  the average per-collective latency, which we can compare against the
  benchmark numbers above.
\end{enumerate}

The key invariant: always measure on the same hardware and software
stack that the production job will use. A benchmark run on a clean VM
with no other traffic will overstate the achievable bandwidth compared
to a crowded cluster.

\section{Common Mistakes}\label{common-mistakes}

\begin{itemize}
\item
  \textbf{Assuming peak bandwidth is sustained.} PCIe Gen5 x16 peaks at
  \textasciitilde64 GB/s, but a single all-reduce never saturates the
  full link because the traffic is split across multiple GPUs and
  multiple links. Measured effective bandwidth is typically 40--60\% of
  peak; design against the measured number, not the spec sheet peak.
\item
  \textbf{Placing all-reduce on Ethernet without a bandwidth budget.} A
  70B model gradient all-reduce at 25 Gb/s takes \textasciitilde56 s per
  step. If the training runs 1,000 iterations, that is
  \textasciitilde56,000 s (15.5 hours) spent exclusively in
  communication. Always compute the per-iteration communication cost
  before committing to a lower-bandwidth fabric.
\item
  \textbf{Ignoring contention from other traffic.} NVSwitch can carry
  many concurrent channels, but if other collectives (e.g., all-gather
  for optimizer state) are running simultaneously, the bisection
  bandwidth is shared. Measure with the full production workload
  pattern, not an isolated micro‑benchmark.
\item
  \textbf{Using InfiniBand as a default without checking the
  cost--performance ratio.} InfiniBand HDR 400 Gb/s gives
  \textasciitilde50 GB/s unidirectional, which is sufficient for many
  workloads, but the price per GB of optics and cables is 5--10× that of
  Ethernet. If the communication time is acceptable on RoCE2, the
  cheaper fabric may be the better architectural choice.
\item
  \textbf{Overlooking protocol overhead.} RoCE and InfiniBand use
  different encoding (8b/10b vs.~128b/130b) and header sizes. A 400 Gb/s
  InfiniBand link delivers \textasciitilde380 Gb/s useful payload; a 200
  Gb/s RoCE link delivers \textasciitilde180 Gb/s. The difference
  matters when we are computing per-second throughput for a fixed data
  volume.
\end{itemize}

\section{Architecture Consequence}\label{architecture-consequence}

The interconnect choice is not a detail to be decided after the model,
the parallelism strategy, and the hardware are selected --- it is a
first-class design variable that shapes all three. If the architecture
calls for data-parallel training of a 70B FP16 model across 8 nodes, the
all-reduce bandwidth dictates whether the nodes are connected by
InfiniBand, 200 GbE, or even 100 GbE. If the budget can accommodate
NVLink‑connected hosts, keeping the all-reduce on-node reduces
communication time by 2 orders of magnitude compared to any multi-node
fabric.

The architecture consequence also flows in the other direction: a
tighter communication budget (e.g., target \textless0.5 s per
all-reduce) may force a particular parallelism strategy. We might choose
model parallelism (pipeline parallelism, tensor parallelism) to keep the
per-collective data volume smaller, or we might accept a slower fabric
and reduce the frequency of collectives (e.g., gradient accumulation
steps between all-reduce calls). The interconnect bandwidth number is
the anchor that makes these trade-offs quantifiable.

Finally, the architecture consequence extends to cost. A cluster of 8 ×
H100 servers connected by NVSwitch is the highest upfront capital
expenditure, but it delivers the lowest per-iteration communication
time, which translates to higher throughput per dollar over the life of
the cluster. A cluster connected by 200 GbE RoCE has lower capital cost
but higher per-iteration communication time, which may increase the
total cost of ownership if the training job runs for many weeks. The
architect must balance capex vs.~opex by way of the communication
arithmetic shown in this chapter.

\section{What We Still Don't Know}\label{what-we-still-dont-know}

\begin{itemize}
\item
  \textbf{Contention patterns under mixed workloads.} The benchmark
  numbers above assume a dedicated collective on a quiet fabric. In a
  production fleet, other traffic (user requests, other training jobs,
  KV-cache movement) shares the same links. The effective bandwidth
  under realistic multi‑job congestion is not yet characterized with a
  standard suite, and the interaction between all-reduce and
  RoCE/Ethernet best‑effort traffic can cause tail latency spikes that
  are difficult to predict.
\item
  \textbf{Adaptive collective algorithms.} New collectives (e.g.,
  gradient compression, sparse all-reduce, error‑bounded approximations)
  claim to reduce the data volume at the cost of a small accuracy loss.
  The bandwidth arithmetic changes when the data volume is 30 GB instead
  of 140 GB, but the latency per byte may increase because of
  compression/decompression steps. The net effect on iteration time is
  workload‑specific and not yet settled.
\item
  \textbf{Optical interconnects and the next generation.} NVLink‑Switch
  and XNAME are pushing node‑level bandwidth higher (toward 1.8 TB/s
  aggregate), and optical intra‑rack fabrics (400 Gb/s--800 Gb/s per
  lane) are entering data centres. The arithmetic base --- the relative
  magnitudes of PCIe, NVLink, Ethernet, InfiniBand --- is stable in the
  near term, but the absolute numbers will shift as new silicon and
  cable technologies ship. Keep the framework; update the constants.
\item
  \textbf{Cross‑domain generalization.} Most published bandwidth
  measurements use square matrices (all-reduce of a gradient tensor).
  Diagonal dominant tensors, sparse gradients, or structured compression
  may exhibit different effective bandwidths because the traffic pattern
  across links is less uniform. The arithmetic framework applies, but
  the constants need re‑measurement for the specific tensor shape.
\end{itemize}

\section{End-of-Chapter Mini-Case}\label{end-of-chapter-mini-case}

An architect is asked to design a distributed inference fleet for a 70B
parameter model serving user queries at 200 tokens/s per user, with
10,000 concurrent users. The team debates whether to use a single 8×H100
host with NVLink, or to shard the model across 32 nodes connected by 200
GbE RoCE. Before any hardware is ordered, the architect opens this
chapter and runs the communication arithmetic.

The per-request all-reduce is not the right model for inference, and
neither is a blanket ``all-gather the full weights every decode step.''
It matters which parallelism is in play: ordinary \textbf{tensor
parallelism (TP)} shards each weight matrix across the GPUs, and every
decode step exchanges only the \emph{activations} (not the weights) via
collective ops like all-reduce/all-gather of the small activation
tensors --- weight matrices stay resident and are read in place, they
are not re-fetched each step. A different scheme, FSDP-style
\textbf{weight gathering} (used in training, and in some
sharded-inference arrangements) does all-gather weight shards before the
forward pass. So the architect must name the scheme before doing the
arithmetic. Concretely, for the canonical 8×H100 host the meaningful
per-step data movement is a TP all-reduce of the per-layer activations
(tens of MBs, not GBs) plus the KV-cache fragment exchange in
disaggregated serving --- e.g.~\textasciitilde35 GB of KV moved over
NVLink (\textasciitilde0.04 s at 900 GB/s) when a P-D sequence transfers
context between prefill and decode pools. The distinction changes the
answer by orders of magnitude: re-fetching 140 GB of weights each step
is pathological (that is a design smell, an anti-pattern to flag, not a
baseline); exchanging activations or KV fragments is the normal case
that fits comfortably within NVLink.

The architect proposes the single 8×H100 host (TP within the node, so
intra-node NVLink carries the per-step activation all-reduce) for the
core serving fleet, with a fallback to a 32‑node RoCE cluster only if KV
concurrency or memory capacity exceeds one host (\hyperref[chap:17]{Chapter 17}). The
communication arithmetic made the trade-off explicit: TP activation
exchanges are tens of MB and free within NVLink; cross-node weight/KV
shuttling is 35× more expensive and avoided unless the fleet genuinely
spans hosts. The design is now grounded in bandwidth and the
\emph{specific} parallelism scheme, not in vague notions of ``scale.''

\begin{center}\rule{0.5\linewidth}{0.5pt}\end{center}
